% ============================================================
%  ch00：标题、摘要、§1 引言（原文 pp.1–4）
% ============================================================
\begin{center}
  {\LARGE\bfseries 几何朗兰兹的一个算术应用}\\[6pt]
  {\large Sam Raskin}\\[2pt]
  {\normalsize Yale University}
\end{center}

\vspace{0.6em}
\begin{center}
\begin{minipage}{0.92\textwidth}
\small
\noindent\textbf{摘要.}\quad Vincent Lafforgue 为函数域上的尖自守函数空间构造了 Langlands 分解。在我们与 Arinkin、Gaitsgory、Kazhdan、Rozenblyum、Varshavsky 的合作工作中，我们证明了几何朗兰兹猜想的一个版本给出 Lafforgue 分解的特征空间在处处无分歧情形下的描述。

在本文中，我们概述后一个思想圈。然后我们说明如何用这些方法证明：在\emph{平凡} Langlands 参数处不存在处处无分歧的紧支撑尖形式，从而回答 Michael Harris 的一个问题。

作为一项具有一定独立意义的结果，我们在平凡 Langlands 参数附近以某些显式项计算伪 Eisenstein 级数的谱类比。在合适的坐标下，我们发现它是 Weyl 特征标公式与一个和曲线相关的 zeta 因子的乘积。
\end{minipage}
\end{center}

\vspace{0.8em}

\section{引言}

\subsection{背景与目标}

\pnum{在我们与 Arinkin、Gaitsgory、Kazhdan、Rozenblyum、Varshavsky 的合作工作 [2]、[3]、[4] 中，我们提出了无分歧、紧支撑自守函数的谱分解的一个猜想性表述。这一猜想也由朱歆文在 [32] 中独立发现。}

我们的谱分解受 V.~Lafforgue 的突破 [22] 与 Beilinson--Drinfeld 的几何朗兰兹猜想启发，属于 Langlands 型的分解，但性质不同：它描述的是全部（紧支撑）自守函数，而不仅是特征函数；并且它同时给出互反律与函子性的陈述，而不必把其中任何一个明确纳入表述。

我们工作的一个主要目的是证明：这一谱分解确实可由几何朗兰兹猜想的 $\ell$-进版本推出。

\pnum{在宣传我们的合作工作时——包括我在 IHÉS 的若干报告——我一直试图论证：我们的猜想能给出关于自守函数的、令数论学家感兴趣的新洞察。我们的猜想在最显而易见的情形下最令人满意，即\emph{离散}（亦称\emph{椭圆}）Langlands 参数（参见例 2.4.5.1）。但在其他 Langlands 参数附近，我一直难以给出精确、具体的推论。}

例如，就目前的形式而言，我们的猜想并不能立即重新得到离散系列的 Arthur 重数公式。要理解这类公式，还需要理论的进一步发展。

\pnum{话虽如此，本文的主要新贡献是给出我们工作对自守函数的一个简单、具体的应用，回答 Michael Harris 的一个问题。这个断言所涉及的是平凡 Langlands 参数——它几乎与离散参数相距最远。}

\pnum{此外，为与会议论文集相称，我们在 §2 中对几何朗兰兹纲领以及 [2]、[3]、[4] 中的思想圈作一些介绍。这两部分可以基本独立地阅读。}

最关心这部分综述材料的读者可以跳到 §2；因为那部分材料按其性质是导引性的，而引言余下部分的重点放在本文后半所考虑的问题上。

\subsection{主结果的陈述}

\pnumt{设定。}{}\ 我们固定一个特征为 $p$ 的有限域 $\F_q$，并令 $k = \Fb$ 表示它的代数闭包。我们固定一个素数 $\ell \neq p$，并令 $\e$ 表示 $\overline{\Q}_\ell$；在 [2] 的术语中，这是\emph{系数域}。我们固定一个分裂约化群 $G/\F_q$，并令 $\bG/\e$ 表示它的朗兰兹对偶群。

令 $X_0/\F_q$ 为一条光滑、射影、几何连通的曲线，并令 $X = X_0 \times_{\F_q} k$ 表示它到 $k$ 的基变换。

我们令 $F = \F_q(X_0)$ 表示与 $X_0$ 相伴的整体域。我们令 $\A$ 表示它的 adele 环，并令 $\OO \subseteq \A$ 表示整 adele 的子环。

我们令 $\Aut^{\unr}_{G,c}$ 表示 $F$ 上处处无分歧、紧支撑自守函数的空间。按定义，这意味着 $\Aut^{\unr}_{G,c}$ 是函数
\[
  G(F)\backslash G(\A)/G(\OO) \longrightarrow \e
\]
的向量空间，且这些函数具有有限支撑。我们令 $\Aut^{\unr}_{\cusp} \subseteq \Aut^{\unr}_{G,c}$ 表示\emph{尖}自守形式的子空间。\footnote{我们提醒读者一个众所周知的事实：在函数域上，尖自守形式\emph{先验地}就是紧支撑的。}

\pnum{我们一次性固定 $X$ 的一个 $k$-点作为我们的基本群的基点；我们在记号中省略它。\footnote{更好的做法是按 $X$ 上 lisse 层的范畴来思考，我们经常这样做。这里用 $\pi_1^{\et}$ 只是为了使讨论中的某些点更具体。}我们令 $\pi_1^{\et}(X)$ 表示 $X$ 的 \et{} 基本群，令 $\pi_1^{\arthm}(X) := \pi_1^{\et}(X_0)$ 表示算术基本群，并令 $\mathcal{W}_X := \pi_1^{\arthm}(X) \times_{\ZZ} \ZZ$ 表示 $X$ 的 Weil 群（赋以其标准拓扑，使得 $\pi_1^{\et}(X) \subseteq \mathcal{W}_X$ 是开子群）。}

\begin{notation}
为确定起见，我们始终使用\emph{几何} Frobenius 约定。于是我们通过一个生成元对应于\emph{几何} Frobenius 元素，来等同 $\widehat{\ZZ} \simeq \pi_1^{\et}(\Spec(\F_q))$。
\end{notation}

\pnumt{Lafforgue--Langlands 分解。}{}\ 暂时假定 $G$ 是半单的，以简化讨论。（本文正文对一般约化群工作。）

一个 \emph{Langlands 参数}是连续同态 $\rho: \mathcal{W}_X \to \bG(\e)$。一个 Langlands 参数称为\emph{半单}的，如果对任意抛物子群 $\check{P} \subseteq \bG$，只要 $\rho$ 通过 $\check{P}(\e)$ 分解，就存在 Levi 因子 $\check{M} \subseteq \check{P}$，使得 $\rho$ 进一步通过 $\check{M}(\e)$ 分解（参见 [2] §3.5--3.6）。

我们现在回忆，[22] 构造了分解
\begin{equation}
  \Aut^{\unr}_{\cusp} \cong \bigoplus_{[\sigma]} \Aut^{\unr}_{\cusp,[\sigma]},
\end{equation}
其中 $[\sigma]$ 取遍半单 Langlands 参数的共轭类。

\begin{remark}
上述讨论对分歧自守函数同样适用。我们的主要结果限制在无分歧情形，因此我们选择在全文强调无分歧设定。
\end{remark}

\pnum{}本文的主要结果如下：

\begin{thmA}
令 $G$ 为半单群（且非平凡群）。令 $\triv: \mathcal{W}_X \to \bG(\e)$ 表示平凡 Langlands 参数，即取值于单位元的常值映射。

假定 [2] 的受限几何朗兰兹猜想及其与 Eisenstein 级的相容性。那么直和项 $\Aut^{\unr}_{\cusp,[\triv]} \subseteq \Aut^{\unr}_{\cusp}$ 为零。换言之，不存在具有平凡 Langlands 参数的无分歧尖形式。
\end{thmA}

\begin{remark}
尽管我们在正文中没有强调这一点，但可以少用一些假设。确切地说，根据 [2] 推论 14.3.5，$\Shv_{\Nilp}(\Bun_G)$ 分解为 $X$ 上半单 $\bG$-局部系统上的一个直和。人们只需要平凡局部系统处的受限几何朗兰兹。我们期待后续工作完全处理这个问题。
\end{remark}

\begin{remark}
粗略地说，论证如下。在 §3 中，我们讨论 Eisenstein 级数在算术朗兰兹的谱侧对应什么。然后在 §4 中，我们在平凡局部系统附近给出 $\LS^{\arthm}_{\bG}$ 的局部坐标（见定理 4.3.3.1），并随后在这些坐标中显式计算谱 Eisenstein 级数（见定理 4.7.2.1）。由此，结果基本上是显然的（见引理 3.7.1.1）。
\end{remark}

\begin{remark}
我们指出，定理 4.7.2.1 证明的一个关键点暗示了某些发散级数的操作与范畴迹方法之间的关系。我们在 §4.6.5 中把关于这个主题的想法——就其目前的形式而言——写出来。这部分材料可以基本独立于本文其余部分阅读。
\end{remark}

\subsection{一些评论}

\pnumt{动机 I。}{}\ $\Aut^{\unr}_{\cusp,[\triv]}$ 的消失是 Beuzart-Plessis--Harris--Thorne 即将开展的工作中的一个组成部分，该工作通过迹公式研究函数域的局部 Langlands 对应。上述定理使他们的结果以几何朗兰兹对应为条件，而后者近年来已取得很大进展。

\pnumt{动机 II。}{}\ 对 $G = PGL_n$，任意尖形式都有不可约的 Langlands 参数；也就是说，在此情形 $\Aut^{\unr}_{\cusp,[\sigma]} = 0$，除非 $\sigma$ 不可约；关于最近的处理，我们参考 [22] 引理 16.4（它沿用了 [20] 的部分内容）；特别参见 [20] 定理 VI.9 的陈述。

然而，对一般的 $G$，情形更复杂：尖形式可以有可约的 Langlands 参数；这与这些尖形式不满足 Ramanujan 猜想有关。

尽管如此，Arthur 的猜想对可能出现的 $\sigma$ 给出了一些限制。首先注意，存在一个典范映射 $\mathcal{W}_X \to \ZZ \xrightarrow{1 \mapsto q} \e^\times$，我们记之为 $\gamma \mapsto |\gamma|$；选取 $\sqrt{q} \in \e^\times$，我们于是得到一个典范映射
\begin{equation}
  \begin{aligned}
    \mathcal{W}_X &\longrightarrow \mathcal{W}_X \times SL_2(\e),\\
    \gamma &\longmapsto \Bigl(\gamma, \begin{pmatrix} \sqrt{|\gamma|} & 0 \\ 0 & 1/\sqrt{|\gamma|} \end{pmatrix}\Bigr).
  \end{aligned}
\end{equation}
Arthur 的猜想预测：$\Aut^{\unr}_{\cusp,[\sigma]}$ 将为零，除非 $\sigma$ 沿 (1.2) 延拓为 $\mathcal{W}_X \times SL_2$ 到 $\bG$ 的一个不可约表示。\footnote{形式上，这意味着我们在 $\e$ 上有一个映射 $SL_2 \to \bG$，以及一个连续映射 $\mathcal{W}_X \to \bG(\e)$，其像与 $SL_2(\e) \to \bG(\e)$ 的像交换。}

\pnum{}假设给定这样一个 Arthur 参数 $\sigma^\sharp: \mathcal{W}_X \times SL_2 \to \bG$，并假设它到 $\pi_1^{\et}(X)$ 的限制是平凡的，于是我们有一个映射 $\sigma^\sharp: \ZZ \times SL_2 \to \bG$。令 $f \in \bg(\e)$ 表示 $\sigma^\sharp\bigl(0, \bigl(\begin{smallmatrix} 1 & 0 \\ 0 & 1 \end{smallmatrix}\bigr)\bigr)$ 的对数，并令 $F \in \bG(\e)$ 表示 $\bigl(1, \bigl(\begin{smallmatrix} \sqrt{q} & 0 \\ 0 & 1/\sqrt{q} \end{smallmatrix}\bigr)\bigr)$ 的像。注意 $F$ 编码了 $\sigma^\sharp$ 的底层 Langlands 参数，并注意 $\Ad_F(f) = qf$。我们还注意，$f$ 必须非零，否则 $\sigma^\sharp$ 不会是不可约的。

于是我们得到：

\begin{conj}
假设 $F \in \bG(\e)$ 是半单元素，并令 $\sigma_F: \mathcal{W}_X \to \bG(\e)$ 表示相应的 Langlands 参数 $\mathcal{W}_X \to \ZZ \xrightarrow{1 \mapsto F} \bG(\e)$。那么 $\Aut^{\unr}_{\cusp,[F]}$ 平凡，除非 $q$ 是 $\Ad_F: \bg \to \bg$ 的一个特征值。
\end{conj}

\begin{remark}
当然，在这个猜想中 $q$ 可以替换为 $q^{-1}$（对 Arthur 的 $SL_2$ 作用 Cartan 对合），这部分反映了该猜想在我们规范化方式的各种修改下（例如几何 Frobenius 与算术 Frobenius）的不变性。
\end{remark}

\pnum{}假定几何朗兰兹，我们将证明与这个猜想相近的结果。确切地说，我们将证明 $\Aut^{\unr}_{\cusp,[\sigma_F]} = 0$，除非 $\Ad_F$ 有一个特征值等于 $H^1_{\et}(X) \times H^2_{\et}(X)$ 中出现的某个 Frobenius 特征值；由于 $q$ 是 $H^2_{\et}(X)$ 上的 Frobenius 特征值，这比猜想的假设略强一些。我们指出，$F$ 为单位元的情形给出定理 A。

\subsection{约定与记号}

我们的希望是，本文能成为进入长论文 [2]、[3]、[4] 的一个入口。我们在 §2 中提供了一些背景。我们也请读者参考 [2] 的引言以获取更多背景。我们还力求在每次使用这些论文中的技术结果时给出精确引用，希望这能帮助读者在这些工作中导航。

我们一般遵循 [2] 的约定与记号；特别参考 §0.9。我们使用\emph{几何域} $k = \Fb$ 来处理曲线 $X$ 的几何、$H$-丛的模叠 $\Bun_H$ 等等，并使用特征 $0$ 的\emph{系数域} $\e = \overline{\Q}_\ell$ 来处理局部系统的几何。在 §2 之外，我们取 $k = \Fb$。$k$ 上的几何是经典代数几何，而 $\e$ 上的代数几何是导出的。我们使用高阶范畴方法。我们的 DG 范畴假定是 $\e$-向量空间上的充实范畴。

对一个在 $k$ 上局部有限型的代数叠 $\mathcal{Y}$，我们令 $\Shv(\mathcal{Y})$ 表示 $\mathcal{Y}$ 上 $\e$-层的 DG 范畴（见 [15] A.1.1 (d$'$)）。我们令 $\qLisse(\mathcal{Y}) \subseteq \Shv(\mathcal{Y})$ 表示\emph{拟 lisse} 复形的子范畴，如 [2] 定义 1.2.6 所述；这些对象的（比如说反常）上同调是 lisse 层的余极限（按通常意义）。

我们一般把 $\infty$-范畴简称为\emph{范畴}，以简化术语。

我们令 $\DGCat_{\mathrm{cont}}$ 表示在连续 DG 函子下余完备（且可及）的 DG 范畴的范畴。我们把 $\DGCat_{\mathrm{cont}}$ 视为配有 Lurie 张量积。我们令 $\Vect \in \DGCat_{\mathrm{cont}}$ 表示 $\e$-向量空间的 DG 范畴，它是幺半结构的单位。

对一个 DG 范畴 $\mathcal{C}$，我们令 $\mathcal{C}^c$ 表示其紧对象子范畴。当 $\mathcal{C}$ 有 $t$-结构时，我们令 $\mathcal{C}^\heartsuit$（相应地 $\mathcal{C}^{\leq 0}$、相应地 $\mathcal{C}^{\geq 0}$）表示 $t$-结构的核心（相应地连通对象子范畴、相应地余连通对象子范畴）。

关于范畴迹方法的背景，我们参考 [15]。

对 $\e$ 上的仿射代数群 $H$，我们提醒：$\LS^{\restr}_H = \LS^{\restr}_H(X)$ 是 $X$ 上（具有受限变差的）$H$-局部系统的模叠。回忆一个 $H$-局部系统不过是一个 $t$-正合（等价地：右 $t$-正合）对称幺半函子 $\Rep(H) \to \qLisse(X)$（等价地：一个对称幺半 $\e$-线性函子 $\Rep(H)^{\heartsuit,c} \to \mathrm{Lisse}(X)^{\heartsuit}$）。因此，我们定义叠 $\LS^{\restr}_H$ 来参数化右 $t$-正合对称幺半 DG 函子 $\Rep(H) \to \qLisse(X)$；更确切地说，$\LS^{\restr}_H$ 的 $S = \Spec(A)$ 点是右 $t$-正合对称幺半函子 $\Rep(H) \to A\text{-}\mathrm{mod}(\qLisse(X))$ 的群胚。

沿几何 Frobenius $\Frob_X: X \to X$ 的拉回定义了一个映射 $\LS^{\restr}_H \to \LS^{\restr}_H$，我们也称之为 \emph{Frobenius}。\footnote{注意这个映射是 $\e$ 上叠的映射。$\LS^{\restr}_H$ 的 Frobenius 可以看作 Frobenius 作用在 $X$ 的（$\e$-向量空间）\et{} 上同调上的非交换/非线性版本。}它的 Frobenius 不动点按定义是叠 $\LS^{\arthm}_H$，它（同义反复地）参数化从 $\Rep(H)$ 到 $X$ 上拟 lisse Weil 层的右 $t$-正合对称幺半函子。

最后，我们始终假定 $p = \mathrm{char}(k)$ 满足 [2] §14.4.1 的（温和）假设。

\subsection{致谢}

我感谢 Michael Harris 提出这个问题。我也感谢 Sasha Braverman 与 Dennis Gaitsgory 对这个主题的兴趣以及有益的谈话。我感谢 Dima Arinkin、Dennis Gaitsgory、David Kazhdan、Nick Rozenblyum 与 Yasha Varshavsky 在这个主题上的合作，感谢 Vincent Lafforgue 与 Cong Xue 的相关讨论。

最后，我感谢 2022 年 IHÉS Langlands 纲领暑期学校的组织者邀请我作报告，并在本文写作期间给予我非凡的耐心。
